%! TEX root = essay.tex
% the main TeX file which is intended to compile, :VimtexReload after adjustment
% :h vimtex-tex-root
\documentclass[../main.tex]{subfiles}
%----------------------------------------------------------------------------------------
%	TITLE SECTION
%----------------------------------------------------------------------------------------
% MAIN TITLE SECTION
\title{
\textbf{Personal Statement} \\
} % Title and subtitle
%\date{\textbf{\DTMtoday}}
\date{\textbf{\today}}
\author{Uliah}

%----------------------------------------------------------------------------------------
% OTHER TITLE SECTION

%\title{\textbf{Sistem Sarana dan Prasarana Jl. Pinggir Laut} \\ {\Large\itshape Infrastructure of Waterfront Parepare City}} % Title and subtitle

%\author{\textbf{Uliah Shafar} \\ \textit{Universitas Diponegoro}} % Author and institution

%\date{\today} % Date, use \date{} for no date

%----------------------------------------------------------------------------------------

\begin{document}
%\maketitle % Print the title section

%----------------------------------------------------------------------------------------
%	ESSAY BODY
%----------------------------------------------------------------------------------------
%\section{Motivations with which you apply for this program}
%\section*{What has motivated you to pursue a graduate degree? Why are you interested in your major?}

\section{Turkiye Perment}

Architecture has been a major influence in my life.
I have studied it maybe almost a decade since I was in a vocational high school until now.
In the beginning to the middle of vocational school, teachers have always been passionate teaching us about building drawing technique, including calculating planning budget and computer-aided drafting.
In the last term of school, teachers allowed me to did internship at the biggest cement company.
The asthonised and diversity experiences that I receive with my peers had became a motivation for me to continue study instead immediately working with the expertise I had had a while years ago.

In contonuing post secondary school, I felt the time was so short. However, it had broaden my percepective on architecture field.
The first time I entered college, it was compulsory in every meeting to finish the drawing and mass study assignment that was unconscionable.
These assignment, in fact, has grown a good work ethic and hard work on me.
Besides all the labour, the lecture sometimes conducted a study trip or just an usual presentation that intrigue me with a digest of architecutre knowledge from all over the world.

In addition to my education in university, I also quite actived on organization. One of foremost organization I have ever involved was PAMIY, which was a group of architectural student from all university in Yogyakata. Its recent activity I followed were constructing a small library in one of the mosque in Yogyakarta suburb called Parangritis and organizing sport competition among architecture student from yogyakarta region.

In contrast to my bachelor's live, I had been drowned more on writing thing in postgraduate rather than involving extracurricular activities. The reason was a covid has disrupted every part of people's activity. However, this has allowed me to finish my master's degree with 4.0 gpa and also increase my writing skill that result a publication. All my writing works were the scope of public space which I care about because it was related to human lives including social, health, and even wellbeing.

Long story short, after graduate from master, I have been accepted in Muhammadiyah Univeristy of Parepare which is a islamic univeristy as a lecture. I have given lectures about three subject in department of city and regional planning. While I have been working there, I realize scarcity of doctorate lecture in university. Of course, this will cause low compentency in teaching and little expert advice on the current city issues regarding its planning. Therefore I intent to persuing doctorate degree that was available out there.

The advanced architecture of turkiye has made it become the best place to go for seeking higher educaion in architectural study. Obviously, architecture in turkiye has long been a role model for worldwide building. For example, we have seen many building nowadays followed turkiye architecture by adding minaret, dome, and even buttresses. What is more is that Turkish architecture successfully refined its identity throughout period.

The city where I was born is similar to tukiye back in the day which need a refinement. Someday, it is new generation that would have higher chance to take part in city's development. The experience with turkiye would allow me to teach different sphere to my student and instill an experience of turkist architecture in my design process. These actions potentially level up my hometown.


%----------------------------------------------------------------------------------------
%	BIBLIOGRAPHY
%----------------------------------------------------------------------------------------

%\bibliographystyle{apalike}

%\bibliography{biblio.bib}



\end{document}
